% ===========================================================================
\chapter{Diagrama Geral de Casos de Uso do Sistema (Visão Macro)}
\label{cap_diagrama_geral_macro}
% ===========================================================================

A visão macro dos casos de uso do SIGEA-GTT fornece uma perspectiva unificada e abrangente de como os diferentes atores interagem com a fronteira do sistema. O modelo comporta vinte e dois casos de uso distribuídos logicamente em seis pacotes funcionais, cobrindo o ciclo completo da auditoria clínica e da gestão pedagógica hospitalar.

% ===========================================================================
\section{Decomposição Modular por Pacotes Funcionais}
\label{sec_decomposicao_pacotes}
% ===========================================================================

Para viabilizar a manutenibilidade, a rastreabilidade e a segregação de responsabilidades de acordo com as normas de Engenharia de Software \cite{sommerville2019}, o ecossistema funcional do SIGEA-GTT está particionado nos seguintes pacotes:

\begin{enumerate}
    \item \textbf{Pacote 1: Segurança e Governança de Acesso}:
    \begin{itemize}
        \item \texttt{UC01}: Autenticar Usuário no Sistema;
        \item \texttt{UC02}: Forçar Redefinição de Senha Provisória;
        \item \texttt{UC03}: Alterar Senha de Acesso pelo Perfil;
        \item \texttt{UC04}: Gerenciar Contas de Usuários.
    \end{itemize}

    \item \textbf{Pacote 2: Parametrização Clínica IHI-GTT e Gravidades}:
    \begin{itemize}
        \item \texttt{UC05}: Parametrizar Módulos e Gatilhos Clínicos GTT;
        \item \texttt{UC06}: Consultar Catálogo Didático de Gatilhos e Gravidades.
    \end{itemize}

    \item \textbf{Pacote 3: Gestão Pedagógica de Turmas e Atividades}:
    \begin{itemize}
        \item \texttt{UC07}: Criar e Configurar Turma Acadêmica;
        \item \texttt{UC08}: Matricular e Desmatricular Discentes na Turma;
        \item \texttt{UC09}: Cadastrar Atividade Prática com Prontuário Simulado;
        \item \texttt{UC10}: Encerrar Período Letivo da Turma.
    \end{itemize}

    \item \textbf{Pacote 4: Execução de Auditoria Clínica e Prontuário Simulado}:
    \begin{itemize}
        \item \texttt{UC11}: Iniciar Sessão de Auditoria Prática;
        \item \texttt{UC12}: Controlar Cronômetro de Auditoria (20 Minutos);
        \item \texttt{UC13}: Rastrear e Vincular Gatilhos no Prontuário;
        \item \texttt{UC18}: Submeter Auditoria Clínica Completa.
    \end{itemize}

    \item \textbf{Pacote 5: Ferramentas da Qualidade e Análise Causal}:
    \begin{itemize}
        \item \texttt{UC14}: Preencher Diagrama de Ishikawa (6M);
        \item \texttt{UC15}: Construir Matriz de Priorização GUT;
        \item \texttt{UC16}: Elaborar Plano de Ação 5W2H;
        \item \texttt{UC17}: Estruturar Ciclo PDCA.
    \end{itemize}

    \item \textbf{Pacote 6: Avaliação Docente e Indicadores Epidemiológicos}:
    \begin{itemize}
        \item \texttt{UC19}: Avaliar Submissão e Emitir Feedback Docente;
        \item \texttt{UC20}: Comparar Gatilhos do Discente com Gabarito Docente;
        \item \texttt{UC21}: Consolidar e Visualizar Indicadores Epidemiológicos IHI;
        \item \texttt{UC22}: Exportar Relatórios de Desempenho e Indicadores Clínicos.
    \end{itemize}
\end{enumerate}

% ===========================================================================
\section{Diagrama Geral UML de Casos de Uso}
\label{sec_diagrama_geral_tikz}
% ===========================================================================

A Figura \ref{fig_diagrama_geral_uc} ilustra o Diagrama Geral de Casos de Uso do SIGEA-GTT, renderizado nativamente em TikZ, enfatizando a fronteira da aplicação, os quatro atores e os principais casos de uso integrados por suas relações de inclusão e extensão.

\begin{figure}[!ht]
\centering
\caption{Diagrama Geral de Casos de Uso do SIGEA-GTT (Visão Macro)}
\label{fig_diagrama_geral_uc}
\resizebox{0.96\textwidth}{!}{%
\begin{tikzpicture}[
    actor/.style={
        draw=clinical-900,
        fill=clinical-100!40,
        thick,
        rounded corners=2pt,
        align=center,
        font=\bfseries\scriptsize,
        minimum width=2.4cm,
        minimum height=1.2cm
    },
    usecase/.style={
        ellipse,
        draw=clinical-700,
        fill=white,
        thick,
        align=center,
        font=\scriptsize,
        minimum width=3.2cm,
        minimum height=1.1cm,
        inner sep=2pt
    },
    usecase_sub/.style={
        ellipse,
        draw=teal!80!black,
        fill=teal!5,
        thick,
        align=center,
        font=\scriptsize,
        minimum width=3.0cm,
        minimum height=1.0cm,
        inner sep=2pt
    },
    link/.style={
        draw=clinical-900!80,
        thick
    },
    include_link/.style={
        draw=clinical-500,
        thick,
        dashed,
        -Latex
    },
    extend_link/.style={
        draw=alert-amber!90!black,
        thick,
        dashed,
        -Latex
    }
]

    % ====================================================
    % FRONTEIRA DO SISTEMA (SYSTEM BOUNDARY)
    % ====================================================
    \draw[rounded corners=8pt, fill=clinical-100!15, draw=clinical-700, line width=1.5pt] 
        (0.5, -14.5) rectangle (14.5, 3.5);
    \node[anchor=north west, font=\bfseries\small\color{clinical-700}] at (0.8, 3.2) 
        {Fronteira do Sistema: SIGEA-GTT};

    % ====================================================
    % ATORES EXTERNOS
    % ====================================================
    % Ator ADMIN (Esquerda Superior)
    \node[actor] (admin) at (-2.2, 1.8) {
        \textbf{Administrador}\\\texttt{\scriptsize (ADMIN)}
    };

    % Ator PROFESSOR (Esquerda Central)
    \node[actor] (prof) at (-2.2, -4.5) {
        \textbf{Docente / Auditor}\\\texttt{\scriptsize (PROFESSOR)}
    };

    % Ator STUDENT (Esquerda Inferior)
    \node[actor] (student) at (-2.2, -10.5) {
        \textbf{Discente / Aluno}\\\texttt{\scriptsize (STUDENT)}
    };

    % Ator SYSTEM_TIMER (Direita Inferior)
    \node[actor, draw=alert-amber!90!black, fill=alert-amber!10] (timer) at (17.2, -10.5) {
        \textbf{Ator Temporal}\\\texttt{\scriptsize (SYSTEM\_TIMER)}
    };

    % ====================================================
    % CASOS DE USO: PACOTE SEGURANÇA E PARAMETRIZAÇÃO
    % ====================================================
    \node[usecase] (uc01) at (3.5, 2.0) {UC01: Autenticar\\no Sistema};
    \node[usecase] (uc02) at (8.0, 2.0) {UC02: Redefinir\\Senha Inicial};
    \node[usecase] (uc04) at (3.5, 0.2) {UC04: Gerenciar\\Usuários};
    \node[usecase] (uc05) at (8.0, 0.2) {UC05: Parametrizar\\Catálogo GTT};
    \node[usecase] (uc06) at (12.2, 0.2) {UC06: Consultar\\Catálogo GTT};

    % ====================================================
    % CASOS DE USO: PACOTE GESTÃO DE TURMAS
    % ====================================================
    \node[usecase] (uc07) at (3.5, -2.5) {UC07: Criar Turma\\Acadêmica};
    \node[usecase] (uc08) at (8.0, -2.5) {UC08: Matricular\\Discentes};
    \node[usecase] (uc09) at (12.2, -2.5) {UC09: Cadastrar\\Atividade Prática};
    \node[usecase] (uc10) at (3.5, -4.5) {UC10: Encerrar\\Turma};

    % ====================================================
    % CASOS DE USO: PACOTE AUDITORIA E PRONTUÁRIO
    % ====================================================
    \node[usecase] (uc11) at (3.5, -7.5) {UC11: Iniciar Auditoria\\Prática};
    \node[usecase] (uc12) at (8.0, -7.5) {UC12: Controlar\\Cronômetro (20m)};
    \node[usecase] (uc13) at (12.2, -7.5) {UC13: Rastrear\\Gatilhos no Prontuário};

    % ====================================================
    % CASOS DE USO: FERRAMENTAS DA QUALIDADE
    % ====================================================
    \node[usecase_sub] (uc14) at (3.5, -10.5) {UC14: Preencher\\Ishikawa 6M};
    \node[usecase_sub] (uc15) at (7.0, -10.5) {UC15: Matriz GUT\\Priorização};
    \node[usecase_sub] (uc16) at (10.5, -10.5) {UC16: Elaborar\\Plano 5W2H};
    \node[usecase_sub] (uc17) at (13.5, -10.5) {UC17: Ciclo PDCA};

    % Submissão Completa
    \node[usecase] (uc18) at (8.0, -12.5) {UC18: Submeter Auditoria\\Clínica Completa};

    % ====================================================
    % CASOS DE USO: AVALIAÇÃO E INDICADORES
    % ====================================================
    \node[usecase] (uc19) at (3.5, -5.8) {UC19: Avaliar Submissão\\e Emitir Nota};
    \node[usecase] (uc20) at (8.0, -5.8) {UC20: Comparar com\\Gabarito Docente};
    \node[usecase] (uc21) at (12.2, -5.8) {UC21: Consolidar\\Indicadores GTT};

    % ====================================================
    % ASSOCIAÇÕES DOS ATORES
    % ====================================================
    % Admin links
    \draw[link] (admin) -- (uc01);
    \draw[link] (admin) -- (uc04);
    \draw[link] (admin) -- (uc05);

    % Professor links
    \draw[link] (prof) -- (uc01);
    \draw[link] (prof) -- (uc06);
    \draw[link] (prof) -- (uc07);
    \draw[link] (prof) -- (uc08);
    \draw[link] (prof) -- (uc09);
    \draw[link] (prof) -- (uc10);
    \draw[link] (prof) -- (uc19);
    \draw[link] (prof) -- (uc20);
    \draw[link] (prof) -- (uc21);

    % Student links
    \draw[link] (student) -- (uc01);
    \draw[link] (student) -- (uc06);
    \draw[link] (student) -- (uc11);
    \draw[link] (student) -- (uc13);
    \draw[link] (student) -- (uc18);

    % Timer links
    \draw[link, draw=alert-amber!90!black, thick] (timer) -- (uc12);

    % ====================================================
    % RELAÇÕES INCLUDE E EXTEND
    % ====================================================
    % UC01 -> UC02 <<extend>> (se must_change_password)
    \draw[extend_link] (uc02) -- node[above, font=\tiny\bfseries] {<<extend>>} (uc01);

    % UC11 -> UC12 <<include>>
    \draw[include_link] (uc11) -- node[above, font=\tiny\bfseries] {<<include>>} (uc12);

    % UC11 -> UC13 <<include>>
    \draw[include_link] (uc11) to[bend left=15] node[above, font=\tiny\bfseries] {<<include>>} (uc13);

    % UC18 -> Ferramentas da Qualidade <<include>>
    \draw[include_link] (uc18) -- node[left, font=\tiny\bfseries] {<<include>>} (uc14);
    \draw[include_link] (uc18) -- node[right, font=\tiny\bfseries] {<<include>>} (uc15);
    \draw[include_link] (uc18) -- node[right, font=\tiny\bfseries] {<<include>>} (uc16);
    \draw[include_link] (uc18) to[bend right=15] node[right, font=\tiny\bfseries] {<<include>>} (uc17);

    % UC19 -> UC20 <<include>>
    \draw[include_link] (uc19) -- node[above, font=\tiny\bfseries] {<<include>>} (uc20);

\end{tikzpicture}
}
\fonte{Elaborado pelo autor (2026).}
\end{figure}

% ===========================================================================
\section{Análise Estrutural das Relações entre Casos de Uso}
\label{sec_analise_relacoes}
% ===========================================================================

Em conformidade com o padrão OMG UML 2.5.1 \cite{omg2017}, a correta aplicação dos estereótipos relacionais assegura precisão semântica e evita ambiguidades arquiteturais:

\begin{enumerate}
    \item \textbf{Relacionamento de Inclusão (\texttt{<<include>>})}:
    \begin{itemize}
        \item Denota uma dependência obrigatória e incondicional onde o caso de uso base incorpora explicitamente o comportamento do caso de uso incluído em um ponto exato do seu fluxo.
        \item No SIGEA-GTT, o caso de uso \texttt{UC11} (Iniciar Sessão de Auditoria Prática) inclui compulsoriamente \texttt{UC12} (Controlar Cronômetro) e \texttt{UC13} (Rastrear Gatilhos no Prontuário). Do mesmo modo, o fechamento da entrega em \texttt{UC18} (Submeter Auditoria Clínica Completa) inclui a validação estrutural das quatro ferramentas da qualidade (\texttt{UC14}, \texttt{UC15}, \texttt{UC16} e \texttt{UC17}), pois a auditoria não pode ser finalizada com artefatos causais vazios.
    \end{itemize}

    \item \textbf{Relacionamento de Extensão (\texttt{<<extend>>})}:
    \begin{itemize}
        \item Modela um comportamento contingencial ou opcional que é executado apenas quando uma condição de extensão específica é atendida em um \textit{Extension Point} definido.
        \item No fluxo de autenticação, \texttt{UC02} (Forçar Redefinição de Senha Inicial) estende \texttt{UC01} (Autenticar Usuário no Sistema) estritamente quando a variável booleana institucional \texttt{must\_change\_password} for verdadeira no registro do usuário. Se o usuário já tiver redefinido sua senha em login anterior, o caso de extensão não é disparado e o fluxo principal prossegue diretamente para o painel de bordo (\textit{dashboard}).
    \end{itemize}

    \item \textbf{Generalização e Especialização de Atores}:
    \begin{itemize}
        \item Embora os três perfis humanos (\texttt{ADMIN}, \texttt{PROFESSOR}, \texttt{STUDENT}) possuam papéis e níveis de autorização distintos no modelo RBAC, todos compartilham a qualidade canônica de usuários autenticados institucionais. Dessa forma, todos executam \texttt{UC01} e herdam a capacidade de alterar sua própria senha voluntariamente em seu perfil (\texttt{UC03}).
    \end{itemize}
\end{enumerate}
